%==============================================================================
% Chapitre 25 - Masse et matériaux des projectiles
%==============================================================================

\section{Introduction}

La masse constitue l'une des caractéristiques fondamentales d'un projectile.

Associée à la vitesse, elle détermine notamment :

\begin{itemize}
    \item l'énergie cinétique ;
    \item la quantité de mouvement ;
    \item la stabilité ;
    \item la sensibilité aux perturbations extérieures.
\end{itemize}

Les matériaux utilisés influencent également de nombreuses propriétés
balistiques.

\begin{aretenir}

Deux projectiles de même calibre peuvent présenter des comportements très
différents en raison de leur masse ou de leur constitution.

\end{aretenir}

\section{La masse}

La masse représente la quantité de matière contenue dans le projectile.

Elle s'exprime dans le Système international en kilogrammes.

Dans le domaine des armes légères, elle est souvent exprimée :

\begin{itemize}
    \item en grammes ;
    \item en grains (\emph{grains} ou \emph{gr}).
\end{itemize}

\begin{exemple}

Un grain correspond à :

\begin{equation}
1~gr = 0,0648~g
\end{equation}

Une balle de 147 grains possède donc une masse d'environ :

\begin{equation}
9,53~g
\end{equation}

\end{exemple}

\section{Influence de la masse}

À vitesse identique, un projectile plus lourd possède :

\begin{itemize}
    \item davantage d'énergie ;
    \item davantage de quantité de mouvement ;
    \item une meilleure inertie.
\end{itemize}

Il tend généralement à mieux conserver sa vitesse.

En revanche :

\begin{itemize}
    \item il nécessite davantage d'énergie pour être accéléré ;
    \item il peut nécessiter un pas de rayure différent.
\end{itemize}

\section{Énergie cinétique}

L'énergie cinétique d'un projectile est donnée par :

\begin{equation}
E_c = \frac{1}{2}mv^2
\end{equation}

où :

\begin{itemize}
    \item $m$ représente la masse ;
    \item $v$ représente la vitesse.
\end{itemize}

Cette énergie constitue une grandeur importante mais ne suffit pas à
caractériser complètement les performances d'un projectile.

\section{Quantité de mouvement}

Une autre grandeur fondamentale est la quantité de mouvement :

\begin{equation}
p = mv
\end{equation}

où :

\begin{itemize}
    \item $p$ représente la quantité de mouvement ;
    \item $m$ la masse ;
    \item $v$ la vitesse.
\end{itemize}

Cette grandeur intervient notamment dans l'étude des chocs et du recul.

\section{Densité}

La densité d'un matériau influence directement la masse qu'il est possible
d'obtenir pour un volume donné.

À géométrie identique :

\begin{itemize}
    \item un matériau dense produit un projectile plus lourd ;
    \item un matériau léger produit un projectile plus léger.
\end{itemize}

\section{Le plomb}

Le plomb a longtemps constitué le matériau principal des projectiles.

Ses principaux avantages sont :

\begin{itemize}
    \item forte densité ;
    \item facilité de mise en forme ;
    \item coût modéré.
\end{itemize}

Ses inconvénients incluent notamment :

\begin{itemize}
    \item sa faible dureté ;
    \item certaines contraintes environnementales.
\end{itemize}

\section{Le cuivre}

Le cuivre et ses alliages sont très largement utilisés.

Ils servent notamment à fabriquer :

\begin{itemize}
    \item les chemises ;
    \item certains projectiles monolithiques.
\end{itemize}

Le cuivre présente :

\begin{itemize}
    \item une bonne résistance mécanique ;
    \item une excellente aptitude à l'usinage ;
    \item une bonne compatibilité avec les canons rayés.
\end{itemize}

\section{L'acier}

L'acier est utilisé dans certaines conceptions spécialisées.

Il permet notamment :

\begin{itemize}
    \item d'augmenter la résistance mécanique ;
    \item d'améliorer certaines capacités de perforation.
\end{itemize}

\section{Le tungstène}

Le tungstène possède une densité particulièrement élevée.

Il est utilisé dans certains projectiles spécialisés nécessitant :

\begin{itemize}
    \item une forte densité ;
    \item un volume réduit ;
    \item une capacité de pénétration importante.
\end{itemize}

\section{Projectiles monolithiques}

Les projectiles monolithiques sont constitués d'un seul matériau.

Ils sont souvent réalisés en :

\begin{itemize}
    \item cuivre ;
    \item alliages de cuivre ;
    \item matériaux spécialisés.
\end{itemize}

Cette approche simplifie la fabrication et offre certaines propriétés
particulières.

\section{Projectiles à plusieurs matériaux}

De nombreux projectiles modernes associent plusieurs matériaux.

On rencontre par exemple :

\begin{itemize}
    \item un noyau ;
    \item une chemise ;
    \item parfois d'autres composants internes.
\end{itemize}

Chaque matériau remplit alors une fonction spécifique.

\section{Répartition des masses}

La masse totale ne suffit pas à décrire complètement un projectile.

La manière dont cette masse est répartie est également importante.

Elle influence notamment :

\begin{itemize}
    \item le centre de gravité ;
    \item les moments d'inertie ;
    \item la stabilité.
\end{itemize}

Cette notion sera développée dans les chapitres suivants.

\section{Masse et stabilité}

La masse agit indirectement sur la stabilité du projectile.

À géométrie identique :

\begin{itemize}
    \item un projectile plus lourd est souvent plus long ;
    \item un projectile plus long nécessite généralement davantage de
    stabilisation.
\end{itemize}

Le lien exact entre masse et stabilité sera étudié ultérieurement.

\section{Masse et simulation}

Dans les modèles balistiques, la masse constitue l'un des paramètres les
plus importants.

Même les modèles les plus simples l'utilisent.

Les modèles avancés prennent également en compte :

\begin{itemize}
    \item les matériaux ;
    \item les densités ;
    \item la répartition interne des masses.
\end{itemize}

\begin{simulation}

Dans un modèle de masse ponctuelle, la masse est souvent le seul paramètre
interne du projectile représenté explicitement.

\end{simulation}

\section{Vers les chapitres suivants}

Après avoir étudié la masse globale du projectile, nous allons nous
intéresser à la répartition de cette masse.

Cette étape nous conduira naturellement vers :

\begin{itemize}
    \item le centre de gravité ;
    \item les moments d'inertie ;
    \item la stabilité.
\end{itemize}

\section{À retenir}

\begin{aretenir}

\begin{itemize}
    \item La masse influence fortement le comportement balistique.
    \item Les matériaux déterminent une grande partie des propriétés du
    projectile.
    \item L'énergie cinétique dépend de la masse et de la vitesse.
    \item La quantité de mouvement constitue une grandeur complémentaire.
    \item La répartition de la masse est aussi importante que la masse totale.
\end{itemize}

\end{aretenir}
